% a06_BALtorus_d5b: continuation of a06_BALtorus_d5 (input from sec/a06_BALtorus_c.tex): the local densities of ln Z^(j), the curvature form, Theorem G, and the eps_0-freedom of the per-step constants.

\begin{lemma}[gauge-invariant local densities of $\ln Z^{(j)}$]\label{lem:a06d:G1}
For $B$ real with (Cur) and $Z$ complex with $w(Z)\le c_*$ define, for $y\in\Lambda$,
\[
Q_y(B+Z)\defeq-\tfrac12\int_0^\infty\Tr_N\big[(1+s)^{-1}\delta_{yy}-((s+X(B+Z))^{-1})(y,y)\big]ds.
\]
Then: \textup{(a)} $F(B)=\sum_{y\in\Lambda}Q_y(B)$, and $\sum_yQ_y(B+Z)=-\frac12\ln\det X(B+Z)$ with the branch continuous from $Z=0$; \textup{(b)} (gauge invariance) $Q_y(B+d\lambda)=Q_y(B)$ for every real $\lambda$ on the fine sites of $\Omega$, $(d\lambda)_b\defeq\xi^{-1}(\lambda(b_+)-\lambda(b_-))$; \textup{(c)} (C-parity) $Q_y(-B)=Q_y(B)$; \textup{(d)} (single locality) for $w(Z),w(Z')\le c_*$, $Q_y(B+Z)$ is analytic and $|Q_y(B+Z)-Q_y(B+Z')|\le C_1\sum_\Delta\hat w_\Delta(Z-Z')e^{-\theta_2d(y,\Delta)}$, $C_1\defeq4NC'_WS^2_{\theta_2}/\gamma_X$, $S_\theta\defeq\sup_z\sum_ye^{-\theta|y-z|}$; \textup{(e)} (multi-point locality) for real $u_1,\dots,u_n$, $n\le\bar n+1$, $|D^nQ_y(B)[u_1,\dots,u_n]|\le C_{(n)}\sum_{\Delta_1,\dots,\Delta_n}\prod_i\hat w_{\Delta_i}(u_i)e^{-(\theta_2/n)\sum_id(y,\Delta_i)}$, $C_{(n)}\defeq2^{n-1}n^n3^de^{\theta_2}C_1c_*^{1-n}$; \textup{(f)} $DQ_y(0)=0$, and for real $Y$ with $sY$ under (Cur) for $s\in[0,1]$ and real $u$, $DQ_y(Y)[u]=\int_0^1D^2Q_y(sY)[Y,u]ds$.
\end{lemma}

\begin{proof}
(a) For a positive symmetric matrix, $\ln X=\int_0^\infty[(1+s)^{-1}-(s+X)^{-1}]ds$ (spectral calculus; the integrand is $O(s^{-2})$); take the trace over $\Lambda\otimes\R^N$, which is $\sum_y\Tr_N$ of the diagonal blocks; for complex $Z$ both sides are analytic on $w(Z)\le c_*$ by \cref{lem:a06d:G0}(ii) and agree on real $Z$. (b) Under $B\to B+d\lambda$, $U(B_b)\to g(b_+)U(B_b)g(b_-)^{-1}$ with $g(x)\defeq\exp(qe_j\lambda(x))$ orthogonal, all commuting; hence $D_B\to gD_Bg^{-1}$, $Q_j(B+d\lambda)=g_{\mathrm{unit}}Q_j(B)g^{-1}_{\mathrm{fine}}$ (the contours start at the unit point), $P(B+d\lambda)=g_{\mathrm{unit}}P(B)g^{-1}_{\mathrm{unit}}$, $G_j\to g_{\mathrm{fine}}G_jg^{-1}_{\mathrm{fine}}$, so $X(B+d\lambda)=g_{\mathrm{unit}}X(B)g^{-1}_{\mathrm{unit}}$ with $g_{\mathrm{unit}}$ block-diagonal over $y$, commuting with the compression to $\Lambda$; therefore $((s+X(B+d\lambda))^{-1})(y,y)=g(y)((s+X(B))^{-1})(y,y)g(y)^{-1}$ and $\Tr_N$ is invariant. (c) Write $q=O(\oplus_i\theta_iJ\oplus0)O^T$ and $c\defeq O(\oplus_i\sigma_z\oplus1)O^T$, orthogonal with $cqc^{-1}=-q$; then $U(-B_b)=cU(B_b)c^{-1}$ and $X(-B)=(1\otimes c)X(B)(1\otimes c)^{-1}$ (for $N=2$, $q=J$, $c=\sigma_z$: complex conjugation). (d) $(s+X')^{-1}-(s+X)^{-1}=(s+X')^{-1}(X-X')(s+X)^{-1}$ with $X\defeq X(B+Z)$, $X'\defeq X(B+Z')$; bound the $(y,y)$ block by \cref{lem:a06d:G0}(ii) for both resolvents and (i) for $X-X'$: $N(2/(s+\gamma_X))^2C'_W\sum_\Delta\hat w_\Delta\sum_{z,z'}e^{-\theta_2|y-z|}e^{-\theta_W(d(z,\Delta)+d(\Delta,z'))}e^{-\theta_2|z'-y|}$; since $\theta_2\le\theta_W/2$ and $|y-z|+d(z,\Delta)\ge d(y,\Delta)$ the double sum is $\le S^2_{\theta_2}e^{-\theta_2d(y,\Delta)}$, and $\int_0^\infty4(s+\gamma_X)^{-2}ds=4/\gamma_X$; analyticity by dominated convergence. (e) Decompose each $u_i=\sum_\Delta u_{i,\Delta}$; fix $(\Delta_1,\dots,\Delta_n)$, $\alpha_i\defeq\hat w_{\Delta_i}(u_i)>0$, and put $g(\tau)\defeq\sum_{S\subset\{1..n\}}(-1)^{n-|S|}Q_y(B+\sum_{i\in S}\tau_iu_{i,\Delta_i})$ on $|\tau_i|\le\rho_i\defeq c_*/(n\alpha_i)$, where $w(\sum\tau_iu_{i,\Delta_i})\le c_*$. For each $i$, pairing $S$ with $S\cup\{i\}$ gives differences with $Z-Z'=\tau_iu_{i,\Delta_i}$, each $\le C_13^de^{\theta_2}(c_*/n)e^{-\theta_2d(y,\Delta_i)}$ by (d); so $|g|\le2^{n-1}3^de^{\theta_2}C_1(c_*/n)\min_ie^{-\theta_2d(y,\Delta_i)}\le\dots e^{-(\theta_2/n)\sum_id(y,\Delta_i)}$. The mixed derivative $\partial_1\cdots\partial_ng(0)=D^nQ_y(B)[u_{1,\Delta_1},\dots,u_{n,\Delta_n}]$; Cauchy on the polydisc divides by $\prod\rho_i$; sum over the $\Delta_i$. (f) By (c), $DQ_y(0)=-DQ_y(0)$; then $\frac d{ds}DQ_y(sY)[u]=D^2Q_y(sY)[Y,u]$, each $sY$ satisfying (Cur).
\end{proof}

\begin{lemma}[curvature form of a background change at one density]\label{lem:a06d:G2}
Let $B,A$ be real bond fields on $\Omega$ such that $B+\sigma A$ satisfies (Cur) for $\sigma\in[0,1]$, $y\in\Lambda$, $R\ge2$ with $C_R(y)\defeq\{x:|x-y|_\infty\le R\}\subset\Omega$. Put $g_B\defeq\norm qe_j\sup\{|(dB)(p)|:p\subset C_R(y)\}$, $g_A$ likewise, $w_B\defeq\norm qe_j\sup_\Omega|B|$, $w_A$ likewise. Then
\[
|Q_y(B+A)-Q_y(B)|\le C_2\big[(g_B+g_A)g_A+e^{-\theta_2R/8}(w_B+w_A+Rg_B+Rg_A)(w_A+Rg_A)\big],\qquad C_2\defeq64C_{(2)}S^2_{\theta_2/4}.
\]
\end{lemma}

\begin{proof}
\emph{Local axial gauge.} For $x\in C_R(y)$ let $\gamma_{y,x}$ be the axis-ordered fine path from $y$ to $x$ (inside the box) and $\lambda_B(x)\defeq\sum_{b\in\gamma_{y,x}}\xi B_b$; extend to $\Omega$ by $\lambda_B(x)\defeq\lambda_B(\pi x)$, $\pi$ the coordinatewise clamp onto $C_R(y)$. $\lambda_B$ is linear in $B$; put $\beta\defeq B-d\lambda_B$, $\alpha\defeq A-d\lambda_A$, so $(B+\sigma A)-d\lambda_{B+\sigma A}=\beta+\sigma\alpha$. \emph{Inside bounds} ($b\subset C_R(y)$): $\xi\beta_b$ is the sum of $\xi^2(dB)(p)$ over the $\le|b-y|_1/\xi+1$ plaquettes of a strip inside $C_R(y)$, so $|\beta_b|\le(|b-y|_1+\xi)\sup_{C_R(y)}|dB|$, likewise for $\alpha$. \emph{Outside bounds}: $\pi$ maps the ends of a bond to one point or to the ends of a bond $b'\subset C_R(y)$, so $|(d\lambda_B)(b)|\le|B_{b'}|+|\beta_{b'}|\le\sup|B|+(3R+1)\sup_{C_R}|dB|$ and $|\beta_b|\le2\sup|B|+(3R+1)\sup_{C_R}|dB|$ for every $b$, likewise for $\alpha$. \emph{Gauge invariance and C-parity}: by \cref{lem:a06d:G1}(b), $Q_y(B+A)-Q_y(B)=Q_y(\beta+\alpha)-Q_y(\beta)=\int_0^1DQ_y(\beta+\sigma\alpha)[\alpha]d\sigma$, and by (f) this is $\int\!\!\int D^2Q_y(s(\beta+\sigma\alpha))[\beta+\sigma\alpha,\alpha]\,ds\,d\sigma$, every background real with (Cur). \emph{Estimate}: apply (e) with $n=2$; call $\Delta$ inside if $\Delta\subset C_{R-2}(y)$ ($d(y,\Delta)\le R$), otherwise $d(y,\Delta)\ge R-3$. Inside, $\hat w_\Delta(\beta+\sigma\alpha)\le(3d(y,\Delta)+5)(g_B+g_A)$ and $\hat w_\Delta(\alpha)\le(3d(y,\Delta)+5)g_A$; outside, $\hat w_\Delta(\beta+\sigma\alpha)\le2(w_B+w_A)+(3R+1)(g_B+g_A)$ and $\hat w_\Delta(\alpha)\le2w_A+(3R+1)g_A$. $\sum_\Delta(3d+5)e^{-(\theta_2/2)d}\le8S_{\theta_2/4}$ and $\sum_{\text{outside}}e^{-(\theta_2/2)d}\le S_{\theta_2/4}e^{-\theta_2(R-3)/4}$; expand the product of the two $\Delta$-sums.
\end{proof}

\begin{lemma}[{Ba{\l}aban's configuration near $\supp\theta_{k+1}$; \cite[pp.~567, 575--577]{Ba2}}]\label{lem:a06d:G3}
Let $0\le j<k$ and $N_k\defeq B^{k-1}(\Lambda^{(k-1)}_2)$ (rescaled to scale $j$), where $\theta_k=1$. \textup{(a)} On $N_k$ the configuration \cite[(2.93)]{Ba2} equals $A^{(k)}$, hence $B\defeq\tilde B+A_1=A^{(k)}-\tilde B^{(k)}$ on $N_k$. \textup{(b)} $\dist_j(\supp\theta_{k+1},N^c_k\cup\Omega^c)\ge D_{kj}\defeq L^{k-j}(r(t_k)-c_b-M)$ in scale-$j$ units.
\end{lemma}

\begin{proof}
(a) $\theta_{l+1}=1$ on $B^l(\Lambda^{(l)}_2)$ \cite[p.~567]{Ba2}, and $B^{k-1}(\Lambda^{(k-1)}_2)\subset B^l(\Lambda^{(l)}_2)$ for $l\le k-1$ by $\Lambda^{(i+1)}_2\subset\Lambda^{(i+1)}_0\subset\Lambda^{(i)'}_7\subset\Lambda^{(i)'}_2$ \cite[pp.~566, 570]{Ba2}; so every term $(1-\theta_{l+1})\theta_lA^{(l)}$, $l\le k-1$, of \cite[(2.93)]{Ba2} vanishes on $N_k$, where $\theta_k=1$; by \cite[(2.89)]{Ba2}, $\theta_{k+1}B^{(k+1)}+\theta_{k+1}A'^{(k)}+\tilde B^{(k)}=\theta_{k+1}A^{(k)}$, so the configuration is $(1-\theta_{k+1})A^{(k)}+\theta_{k+1}A^{(k)}=A^{(k)}$. (b) $\supp\theta_{k+1}$ lies in the $M$-neighbourhood of $B^k(\Lambda^{(k)}_2)$; $\Lambda^{(k)}_2\subset\Lambda^{(k-1)'}_7$, which lies five layers of width $\ge r(t_{k-1})-c_b$ (scale $k-1$) inside $\Lambda^{(k-1)'}_2$, i.e.\ $\ge r(t_k)-c_b$ in scale-$k$ units at $L=5$; multiply by $L^{k-j}$.
\end{proof}

The inputs read from \cite{Ba2}: (R1) $|\partial^\eta A^{(k)}|\le C_{Z2}p(t_k)$ on $B^k(\Lambda^{(k-1)'}_2)\supset N_k$ \cite[(2.60)]{Ba2}; (R2) $|\tilde B^{(k)}|\le C_\kappa t_k^\kappa$ and $|\partial^\eta\tilde B^{(k)}|\le C'_\kappa t_k^\kappa$ for every $\kappa$ \cite[(2.90)]{Ba2}; (R3) (Cur) at scale $j$ on $\Omega$ for $\tilde B+A_1+\sigma A_2$, $\sigma\in[0,1]$: by \cite[(2.94), (2.98)]{Ba2} the curvature amplitude is $\le4\norm qC_{98}et_j^{1/2}p(t_j)$, which is $\le c_{\mathrm{cur}}$ under the cap $\norm qC_{98}et^{1/2}p(t)\le c_{\mathrm{cur}}/4$ for $t\le\eps_0$; (R4) $w_B\le C_wt_j^{-1/2}p(t_j)$ on $\Omega$, from the restrictions \cite[(2.55)]{Ba2} at each scale $l\in(j,k+1]$ with \cite[(2.59)]{Ba2} and the rescaling $e_jA^{(l),(j)}=L^{-(l-j)}e_lA^{(l)}$.

\begin{theorem}[the terms containing $\tilde B^{(k)}$ are summable over $j<k$]\label{thm:a06d:G}
Let $1\le k\le K-1$ and $j<k$. Then for every $\kappa$,
\[
|D_j|\le C_G(\kappa)(1+(k-j))^{4+\bar p}L^{-(k-j)}t_k^{1+\kappa}p(t_k)|\Lambda^{(k)}_2|,\quad\text{hence}\quad\sum_{j<k}|D_j|\le C'_G(\kappa)t_k^\kappa|\Lambda^{(k)}_2|\ \text{ and }\ |D_j|\le C''_G(\kappa,\beta)t_k^\kappa|\Lambda^{(k)}_2|L^{(j-k)(1-\beta)}
\]
for every $\beta>0$, with $C_G,C'_G,C''_G$ depending only on $(d,N,\norm q,a,L,M,R,\varrho,\bar p,b_0,e,\lmu,c_1,C_{Z2},C_\kappa,C'_\kappa,C_{98},C_w)$ and the constants of \cref{lem:a06d:G1}, free of $j,k,\eps,\eps_0$, the torus and the history.
\end{theorem}

\begin{proof}
Fix $j<k$, $A\defeq A_2$, $B\defeq\tilde B+A_1$ at scale $j$, and $R\defeq R_{kj}\defeq4+(32/\theta_2)(k-j)\ln L$; require $R\ge R_*(\theta_2,L,c_b,M)$, fixed before $\eps_0$, so that $2R_{kj}+3\le D_{kj}$ for all $k-j\ge1$ (possible since $D_{kj}\ge L^{k-j}R/2$ grows faster than $R_{kj}$). \emph{Count}: $\supp A\subset\supp\theta_{k+1}$, and the set $S_A$ of scale-$j$ unit cubes meeting it has $|S_A|\le(2M+1)^3L^{3(k-j)}|\Lambda^{(k)}_2|$. \emph{Class I}, $y$ with $d(y,S_A)\le R$: by \cref{lem:a06d:G3}(b), $C_R(y)\subset\Omega\cap N_k$, and by (a) $B=A^{(k)}-\tilde B^{(k)}$ there, so by (R1), (R2) and the rescaling, $g_B\le2\norm qeL^{-2(k-j)}t_k^{1/2}(C_{Z2}p(t_k)+C'_\kappa t_k^\kappa)$, $g_A\le2\norm qeL^{-2(k-j)}t_k^{1/2}C'_\kappa t_k^\kappa$, $w_A\le\norm qeL^{-(k-j)}t_k^{1/2}C_\kappa t_k^\kappa$, and $w_B\le C_wt_j^{-1/2}p(t_j)=C_wL^{(k-j)/2}t_k^{-1/2}p(t_j)$ with $p(t_j)\le(1+(k-j)\ln L)^{\bar p}p(t_k)$ (R4). \cref{lem:a06d:G2} (hypotheses by (R3)): the main term is $\le CL^{-4(k-j)}t_k^{1+\kappa}p(t_k)$; the other is $\le Ce^{-\theta_2R/8}R^2L^{(k-j)/2}L^{-(k-j)}(1+(k-j)\ln L)^{\bar p}t_k^\kappa p(t_k)\le CR^2(1+(k-j)\ln L)^{\bar p}L^{-4.5(k-j)}t_k^\kappa p(t_k)$, since $e^{-\theta_2R/8}\le L^{-4(k-j)}$. The number of class-I $y$ is $\le(2R+3)^3|S_A|$. \emph{Class II}, $d(y,S_A)>R$: by \cref{lem:a06d:G1}(d) with $Z=A$, $Z'=0$ ($w(A)\le c_*$ for $t_k\le\eps_0$ by (R2), a cap), $|Q_y(B+A)-Q_y(B)|\le C_1w_A\sum_{\Delta\in S_A}e^{-\theta_2d(y,\Delta)}$, which summed over class II is $\le C_1w_A|S_A|S_{\theta_2/2}e^{-\theta_2R/2}\le CL^{3(k-j)}L^{-(k-j)}L^{-16(k-j)}t_k^{1/2+\kappa}(2M+1)^3|\Lambda^{(k)}_2|$. \emph{Sum}: by \cref{lem:a06d:G1}(a), $D_j=\sum_{y\in\Lambda}[Q_y(B+A)-Q_y(B)]$; class I gives $(2R+3)^3(2M+1)^3L^{3(k-j)}|\Lambda^{(k)}_2|\times[CL^{-4(k-j)}t_k^{1+\kappa}p+CR^2(1+(k-j)\ln L)^{\bar p}L^{-4.5(k-j)}t_k^\kappa p]$, whose first part is $\le C(1+(k-j))^3L^{-(k-j)}t_k^{1+\kappa}p|\Lambda^{(k)}_2|$ and whose second part is smaller by $L^{-1.5(k-j)}$ up to polynomials at the cost of one power of $t_k$ (take $\kappa$ one larger in (R2)); class II is smaller still. \emph{Sums over $j$}: $\sum_{m\ge1}(1+m)^{4+\bar p}L^{-m}<\infty$ and $t^{1+\kappa}p(t)\le P(1,1)t^\kappa$; for the per-$j$ form absorb $(1+m)^{4+\bar p}$ into $L^{\beta m}$.
\end{proof}

Per scale-$j$ cube the bound is $\sim L^{-4(k-j)}t_k^{1+\kappa}p=t_j^4t_k^{\kappa-3}p$, i.e.\ $O(t_j^{4-\beta})$ for $\kappa\ge3$: the printed claim ``$O((L^j\eps)^\kappa)|B^{k-j}(\Lambda^{(k)}_2)|$ with $\kappa>d$, $\kappa_0>0$'' holds with $\kappa=4-\beta$, $\kappa_0=1-\beta$, the double-curvature count.

\begin{corollary}[the $Z^{(j)}$ expansions, consumed form]\label{cor:a06d:G}
For $1\le k\le K-1$, $\sum_{j<k}|F_j(\tilde B+A_1+A_2)-F_j(\tilde B)-P_{\tilde B}[A_1]|\le(C_\kappa/(L^{\bar n+1-d}-1)+C'_G(\kappa))t_k^\kappa|\Lambda^{(k)}_2|$ for every $\kappa<(\bar n+1)/2$, with per-$j$ form $O(t_k^\kappa)|\Lambda^{(k)}_2|L^{(j-k)\kappa_0}$, $\kappa_0=\min(\bar n+1-d,1-\beta)>0$. This covers the three categories of the sentence of \cite[p.~578]{Ba2}: the order-$(\bar n+1)$ trace term, the R-vertices and the orders $>\bar n$ ($R_j$), and the terms containing $\tilde B^{(k)}$ ($D_j$); \cite[(2.104)]{Ba2} is unchanged.
\end{corollary}

Every gauge step above is the exact invariance of a full object, $Q_y$ or $F_j$, under a gauge transformation of the background on all of $\Omega$; the local axial gauge of \cref{lem:a06d:G2} is used only to bound fields, never to split an identity, and no Ward identity on a photon line enters.

\subsubsection*{The per-step large-field constants are free of $\eps_0$}

The order of constants in \cite[Sect.~3]{Ba2} requires that each per-step constant of the large-field chain be fixed before $\eps_0$, so that the final caps on $\eps_0$ (\cref{lem:s14:CW}) are consistent. A constant can depend on $\eps_0$ only through a cap ``for $e(t)$ sufficiently small'' or through a bound stated at the final scale; in the first case the constant is the sup over $t\in(0,1]$ of an explicit function, finite by $E(\kappa,\delta)$, $P(a,m)<\infty$, in the second case it is replaced by the scale-local bound. The constants of the chain are: $c_{3.29}$ of \cite[(3.29)]{Ba2}, from \cite[Prop.~3.1$'$, Sect.~4]{Ba4}, whose remainder $O(1)e^2p^2(e)(|\phi(x)|^2+|\phi(x')|^2)$ in \cite[(4.7)]{Ba4} is built from \cite[Lemma~2.1, (4.10), (4.13), (4.22)]{Ba4} with $\gamma''_0$ depending ``on $d$ and $a$ only'', so that $c_{3.29}(k)=C_{\mathrm{Ba4}}e^{2-\alpha}C_\phi^2\lambda^{-1/2}t_k^{(1-\alpha)/2}p_k^2$ is scale-local; $C_{2.57}$ of \cite[(2.57)]{Ba2}, which depends on $(d,L,M,\alpha_0,\bar n,a,c_1,b_0,\bar p,R,\varrho,e,\lambda,N,\beta)$ and on the $O(1)$ of Proposition~1, with $\eps_0$ entering only through the caps $p(t)^n\lambda(t)^{-m}e^{-\delta r(t)}\le t^\kappa$ (Ba{\l}aban's order of choices, \cite[p.~421]{Ba3}); the counterterm constants, $|\dm|\le c_m\eps^{-1}$ (per-graph power counting, no $\eps_0$) and the extensive bounds on the vacuum energy and on $\dm_{K,k}$ (Step~2 and \cref{thm:a06d:V12,thm:a06d:V14}), whose constants are built from renormalised power counting; and the $Z$-factor constant $C'_G$ of \cref{thm:a06d:G}, whose caps are preliminary. So every per-step constant is free of $\eps_0$, and the order statement of \cref{thm:a06:BALA} holds.
